\section*{7. 有限群不变量的有限性定理}
\begin{center}
Math. Ann. 77 (1916), pp. 89--92
\end{center}

% unit: noether.p07.par.001
以下将给出一个关于有限群不变量有限性的完全初等的证明；它只依赖于对称函数理论，同时还给出完整系统的实际指定。通常的证明则依靠 Hilbert 的模基定理（Ann. 36），因而只是存在性证明。\footnote{参见例如 Weber, \emph{Lehrbuch der Algebra} (2nd ed.), vol. 2, \S\ 57.}

% unit: noether.p07.par.002
设有限群 $\Hgrp$ 由 $h$ 个线性变换（行列式不为零）$A_1,\ldots,A_h$ 构成，其中 $A_k$ 表示线性变换
% unit: noether.p07.eq.001
\[
  x_1^{(k)}=\sum_{\nu=1}^{n}a_{1\nu}^{(k)}x_\nu,\ldots,
  x_n^{(k)}=\sum_{\nu=1}^{n}a_{n\nu}^{(k)}x_\nu,
\]
或简写为 $(x^{(k)})=A_k(x)$。因此，群 $\Hgrp$ 把具有元素 $x_1,\ldots,x_n$ 的行 $(x)$ 变为具有元素 $x_1^{(k)},\ldots,x_n^{(k)}$ 的各行 $(x^{(k)})$。由于恒等变换必须包含在 $A_1,\ldots,A_h$ 之中，行 $(x)$ 本身也包含在这些行 $(x^{(k)})$ 之中。所谓该群的整有理（绝对）不变量，是指 $x_1,\ldots,x_n$ 的这样的整有理函数：在施用 $A_1,\ldots,A_h$ 时恒等地保持不变。因此，对于这样的不变量 $f(x)$，有
% unit: noether.p07.eq.002
\begin{equation}
  f(x)=f(x^{(1)})=\cdots=f(x^{(h)})=
  \frac1h\sum_{k=1}^{h}f(x^{(k)}).
\tag{1}
\end{equation}

% unit: noether.p07.sec.001
\subsection*{1.}
% unit: noether.p07.par.003
公式 (1) 表明，$f(x)$ 是诸数量行 $(x^{(1)}),\ldots,(x^{(h)})$ 的整有理对称函数；而且，由于每一项 $f(x^{(k)})$ 只含有一个行 $(x^{(k)})$，这里涉及的是最简单的情形，通常称为单式情形。由关于数量行的对称函数的已知定理，\footnote{参见 2. 的注。}$f(x)$ 因而能由这些行的初等对称函数整有理地表示，也就是说，由“Galois 预解式”的系数 $G_{\alpha\alpha_1\ldots\alpha_n}(x)$ 表示：
% unit: noether.p07.eq.003
\[
\begin{aligned}
  \Phi(z,u)
  &=\prod_{k=1}^{h}\bigl(z+u_1x_1^{(k)}+\cdots+u_nx_n^{(k)}\bigr) \\
  &=z^h+
    \sum G_{\alpha\alpha_1\ldots\alpha_n}(x)
    z^{\alpha}u_1^{\alpha_1}\cdots u_n^{\alpha_n},
\end{aligned}
\qquad
\left(\begin{array}{c}
\alpha+\alpha_1+
\cdots+
\alpha_n=h\\
\alpha\ne h
\end{array}\right),
\]
其中 $G_{\alpha\alpha_1\ldots\alpha_n}(x)$ 是关于 $x$ 的次数为 $\alpha_1+\cdots+\alpha_n$ 的不变量。于是得到证明：

% unit: noether.p07.thm.001
\emph{Galois 预解式的系数 $G_{\alpha\alpha_1\ldots\alpha_n}(x)$ 构成该群的不变量完整系统，使每一个不变量都能由这有限多个不变量整有理地表示。}

% unit: noether.p07.sec.002
\subsection*{2.}
% unit: noether.p07.par.004
从 (1) 出发，还可以作如下更为初等的考察，\footnote{这等同于在上述单式情形中证明关于数量行对称函数的定理。}它同时导向第二个完整系统。设
% unit: noether.p07.eq.004
\[
  f(x)=a+b x_1^{\mu_1}\cdots x_n^{\mu_n}
       +\cdots+c x_1^{\nu_1}\cdots x_n^{\nu_n},
\]
其中 $a,b,\ldots,c$ 表示常数。由 (1) 得
% unit: noether.p07.eq.005
\[
\begin{aligned}
 h\cdot f(x)
 &=h\cdot a+b\sum_{k=1}^{h}(x_1^{(k)})^{\mu_1}\cdots(x_n^{(k)})^{\mu_n}
     +\cdots
     +c\sum_{k=1}^{h}(x_1^{(k)})^{\nu_1}\cdots(x_n^{(k)})^{\nu_n} \\
 &=h\cdot a+b\cdot J_{\mu_1\ldots\mu_n}+\cdots+c\cdot J_{\nu_1\ldots\nu_n}.
\end{aligned}
\]
因此，每个不变量都可以由特殊不变量
% unit: noether.p07.eq.006
\[
  J_{\mu_1\ldots\mu_n}
  =\sum_{k=1}^{h}(x_1^{(k)})^{\mu_1}\cdots(x_n^{(k)})^{\mu_n}
\]
整线性地合成，所以只需对这些特殊不变量证明断言。但除去一个数值因子外，$J_{\mu_1\ldots\mu_n}$ 是表达式
% unit: noether.p07.eq.007
\[
  S_\mu=\sum_{k=1}^{h}
  \bigl(u_1x_1^{(k)}+\cdots+u_nx_n^{(k)}\bigr)^{\mu}
\]
中 $u_1^{\mu_1}\cdots u_n^{\mu_n}$ 的系数，其中 $\mu=\mu_1+\cdots+\mu_n$；而该表达式是 $h$ 个线性形式
% unit: noether.p07.eq.008
\[
  \xi_1=u_1x_1^{(1)}+\cdots+u_nx_n^{(1)},\ldots,
  \xi_h=u_1x_1^{(h)}+\cdots+u_nx_n^{(h)}
\]
的第 $\mu$ 个幂和。众所周知，无穷多个幂和 $S_\mu$ 都是
% unit: noether.p07.eq.009
\[
  S_1,S_2,\ldots,S_h
\]
的整有理组合，其系数由不变量
% unit: noether.p07.eq.010
\[
  J_{\mu_1\ldots\mu_n}\qquad(\mu_1+\cdots+\mu_n\le h)
\]
给出。于是得到第二个完整系统：

% unit: noether.p07.thm.002
\emph{该群的不变量完整系统由所有满足 $\mu_1+\cdots+\mu_n\le h$ 的不变量 $J_{\mu_1\ldots\mu_n}$ 给出，其中 $h$ 表示群的阶。}

% unit: noether.p07.par.005
与 1. 的联系由如下观察建立：幂和 $S_1,\ldots,S_h$ 可以用 $\xi_1,\ldots,\xi_h$ 的初等对称函数替代，这就导向那里给出的 $G_{\alpha\alpha_1\ldots\alpha_n}(x)$ 的完整系统。两个结果都表明，所有不变量都可以由那些关于 $x$ 的次数不超过群阶 $h$ 的不变量整有理地表示。

% unit: noether.p07.sec.003
\subsection*{3.}
% unit: noether.p07.par.006
由刚刚证明的结果可以推出关于有理表示的结论。每一个有理绝对不变量都可以表示为两个整有理绝对不变量的商，这两个不变量不必互素；这是通过把分子和分母乘以分母关于该群的共轭即可直接看出的。因此，每一个有理绝对不变量都可以由 Galois 预解式 $\Phi(z,u)$ 的系数 $G_{\alpha\alpha_1\ldots\alpha_n}(x)$ 有理地表示。这个定理也可以不经由有限性定理而以多种方式容易地证明；它在 Weber II, \S\ 58 中已经表述。\footnote{然而，那里给出的证明并不可靠。它只表明公式 (7) 中的函数 $\Psi(t)$ 以不变量为系数，却没有表明这些不变量能由 $\Phi(t)$ 的系数有理表示。若在用 $\xi_k$ 表示 $x_i^{(k)}$ 时，不使用 Lagrange 插值公式，而使用一个已知的微分方法，则可避免这个缺口。这就是对恒等式 $\Phi(-\xi_k,u)=0$ 关于所有 $u$ 微分得到的关系：
\[
 \left[\frac{\partial\Phi}{\partial u_i}
      -x_i^{(k)}\frac{\partial\Phi}{\partial z}\right]_{z=-\xi_k}=0.
\]
相应地，对于每一个有理函数 $\omega(x)$，Weber 中的公式 (7) 和 (8) 应由下式代替：
\[
 \omega(x_1^{(k)},\ldots,x_n^{(k)})=
 \left.
 \omega\!\left(
   \frac{\partial\Phi/\partial u_1}{\partial\Phi/\partial z},\ldots,
   \frac{\partial\Phi/\partial u_n}{\partial\Phi/\partial z}
 \right)\right|_{z=-\xi_k},
\]
它只含有 $\Phi(z,u)$ 的系数；当 $\omega(x)$ 为不变量时，对所有 $k$ 求和即给出所要求的表示。}

% unit: noether.p07.sec.004
\subsection*{4.}
% unit: noether.p07.par.007
最后要指出，这里推出的结果已隐含于我的论文“有理函数的域与系统”之中；\footnote{Math. Ann. 76, p. 161 (1915).}更准确地说，1. 中给出的完整系统包含在定理 VI 和 VII 中，而独立于这个有限性定理的有理表示证明包含在定理 III 中。\footnote{对于相对不变量，定理 VIII 和 IX 包含了一个有限性证明，它同样建立在不同于通常证明的基础上，但也只是存在性证明；原因在于相对不变量并不构成一个域。}不过，在目前这个特殊情形中，证明过程和结果相对于一般理论都大为简化。我之所以想到把一般研究应用于有限群不变量，是由于与 E. Fischer 的谈话；2. 中给出的证明在内容上也来自他，3. 的注释亦然；在一般情形中，那里出现的完整系统并不能有理定义。

% unit: noether.p07.close.001
Erlangen, 1915 年 5 月。

\clearpage
% unit: noether.p08.title.001
